\documentclass[11pt]{amsart}
\usepackage[T1]{fontenc}\usepackage{lmodern}
\usepackage{amsmath,amssymb,amsthm,mathtools,xcolor}
\usepackage[colorlinks=true,linkcolor=blue!55!black,citecolor=blue!55!black,urlcolor=blue!55!black]{hyperref}
\newtheorem{theorem}{Theorem}[section]\newtheorem{lemma}[theorem]{Lemma}
\newtheorem{proposition}[theorem]{Proposition}\theoremstyle{definition}
\newtheorem{remark}[theorem]{Remark}
\title[A Preorder Support-Polynomial Counterexample]{An Eight-Element Counterexample to Real-Rootedness of Preorder Support Polynomials}
\author{Zijian Zeng}
\address{Institute of Computer Science and Digital Innovation, UCSI University, Kuala Lumpur, 56000, MALAYSIA}
\email{zijianzeng@foxmail.com}\email{1002266693@ucsiuniversity.edu.my}
\date{August 31, 2026}
\hypersetup{pdftitle={An Eight-Element Counterexample to Real-Rootedness of Preorder Support Polynomials},pdfauthor={Zijian Zeng}}
\begin{document}
\begin{abstract}
We independently certify a recently posted eight-element counterexample to the
real-rootedness conjecture for the support polynomial of a preorder polytope.
The witness is an ordinary height-two poset. Its support polynomial is computed
from all order-ideal inequalities, and non-real-rootedness is proved by a
rational Laguerre-inequality certificate, with an independent Sturm-chain
check. Thus the computational frontier at size eight is negative.
\end{abstract}\maketitle

\section{Introduction}
Athanasiadis and Chapoton introduced preorder polytopes and conjectured that
their lattice-point support polynomials have only real roots
\cite[Conjecture~5.3]{chapoton2026polytopesposetsassociatedpreorders}. They reported computer verification for every
preorder of size at most seven and every arbor of size at most ten
\cite[p.~16]{chapoton2026polytopesposetsassociatedpreorders}. This support polynomial is not the Ehrhart
$h^*$-polynomial.

A candidate eight-element counterexample was subsequently posted publicly on
MathDB \cite{Patel2026}. This note gives a self-contained exact verification
from the defining inequalities, together with two algebraic certificates of
non-real-rootedness. A later work of Dai, Hou, Liu, Thawinrak and Wang proves
the separate duality conjecture for preorder support polynomials, not the
real-rootedness conjecture \cite{dai2026countinglatticepointsminkowski}.

\section{Definitions and the witness}
Let $P$ be a finite poset. An order ideal $I\subseteq P$ satisfies $x\in I$
and $y\le x$ only if $y\in I$. Following \cite{chapoton2026polytopesposetsassociatedpreorders}, set
\begin{align*}
 Q_P&=\left\{x\in\mathbb R_{\ge0}^{P}:\sum_{v\in I}x_v\le |I|
 \text{ for every order ideal }I\right\},\\
 h_P(t)&=\sum_{x\in Q_P\cap\mathbb Z^P}t^{|\{v:x_v>0\}|}.
\end{align*}
Consider the height-two poset with lower elements $a_1,a_2,a_3,d$, upper
elements $b_0,b_1,b_2,b_3$, and strict relations
\begin{equation}\label{eq:relations}
 a_i<b_0,\qquad a_i<b_i,\qquad d<b_i\qquad(1\le i\le3).
\end{equation}
There are no other strict relations. This is an ordinary poset, hence a
preorder of size eight.

\section{Exact support enumeration}
\begin{proposition}\label{prop:h}
For the poset in \eqref{eq:relations},
\[
h_P(t)=1+17t+106t^2+303t^3+427t^4+303t^5+106t^6+17t^7+t^8.
\]
\end{proposition}
\begin{proof}
Each lower coordinate is $0$ or $1$, because its singleton is an order ideal.
Fix the subset $Z$ of lower elements having coordinate zero, and put
\[r=|Z\cap\{a_1,a_2,a_3\}|,\qquad
\epsilon=\boldsymbol1_{d\in Z}.\]
There are $\binom3r$ choices of this type. By symmetry take the zero $a_i$'s
to be $a_1,\ldots,a_r$. Write $y_i=x_{b_i}$ and
$e_i=\max(y_i-1,0)$. For an upper subset $S$, the least ideal containing it is
$S\cup N(S)$, where $N(S)$ is its lower neighbourhood. Any larger ideal only
adds lower coordinates $x_v\le1$ and gives no stronger inequality.
Consequently the order-ideal inequalities are equivalent to
\[
0\le e_0\le r,\qquad
\sum_{i=1}^3(e_i-\boldsymbol1_{i\le r})_+\le\epsilon,\qquad
e_0+e_1+e_2+e_3\le r+\epsilon.
\]
Let $F_{r,\epsilon}(u)$ enumerate these vectors by excess support. Bounded
counting gives
\[
\begin{array}{c|c@{\quad}c|c}
(r,\epsilon)&F_{r,\epsilon}(u)&(r,\epsilon)&F_{r,\epsilon}(u)\\ \hline
(0,0)&1 &(0,1)&1+3u\\
(1,0)&1+2u &(1,1)&1+5u+5u^2\\
(2,0)&1+4u+3u^2 &(2,1)&1+7u+13u^2+4u^3\\
(3,0)&1+6u+9u^2+4u^3 &(3,1)&1+9u+24u^2+16u^3+u^4.
\end{array}
\]
For an excess vector with $s$ positive coordinates, every zero-excess upper
coordinate is independently $0$ or $1$. Hence
\[
h_P(t)=\sum_{r=0}^3\sum_{\epsilon=0}^1\binom3r
t^{4-r-\epsilon}(1+t)^4F_{r,\epsilon}\!\left(\frac{t}{1+t}\right).
\]
Expansion proves the formula. Separately, enumeration of all 44 order ideals
and bounded integer vectors gives the same coefficients and $h_P(1)=1281$.
\end{proof}

\section{Non-real-root certificate}
\begin{lemma}\label{lem:laguerre}
If a real polynomial $q$ has only real roots, then
$(q'(s))^2-q(s)q''(s)\ge0$ for every real $s$.
\end{lemma}
\begin{proof}
Away from the roots, logarithmic differentiation of
$q(s)=c\prod_j(s-\rho_j)$ gives
\[(q')^2-qq''=q(s)^2\sum_j(s-\rho_j)^{-2}\ge0.\]
At roots the result follows by continuity; multiplicities cause no change.
\end{proof}
\begin{theorem}\label{thm:main}
The polynomial $h_P(t)$ in Proposition~\ref{prop:h} has non-real roots.
Consequently Conjecture~5.3 of \cite{chapoton2026polytopesposetsassociatedpreorders} is false, already for an ordinary
poset of size eight.
\end{theorem}
\begin{proof}
Coefficient comparison gives
\begin{equation}\label{eq:gamma}
h_P(t)=(1+t)^8G\!\left(\frac{t}{(1+t)^2}\right),\quad
G(z)=1+9z+24z^2+16z^3+z^4.
\end{equation}
At $s=-1/4$,
\[G(s)=\frac1{256},\quad G'(s)=-\frac1{16},\quad G''(s)=\frac{99}{4},\]
and therefore $G'(s)^2-G(s)G''(s)=-95/1024<0$.
Lemma~\ref{lem:laguerre} implies that $G$ has a non-real root $z$. Let $t$
be either root of $z(1+t)^2-t=0$. Such a $t$ is not $-1$, and
\eqref{eq:gamma} gives $h_P(t)=0$. If $t$ were real, then
$z=t/(1+t)^2$ would be real, a contradiction.
\end{proof}
\begin{remark}[Independent Sturm check]
A Sturm chain for $G$ is
\[
\begin{aligned}
S_0&=z^4+16z^3+24z^2+9z+1,&S_1&=4z^3+48z^2+48z+9,\\
S_2&=144z^2+165z+32,&S_3&=651z+160,&S_4&=-1.
\end{aligned}
\]
The successive negative remainders are $1/4$, $7/1728$ and $6848/47089$
times the next polynomial. There are three sign variations at $-\infty$ and
one at $+\infty$, so $G$ has exactly two real roots and one non-real pair.
\end{remark}

\section*{Data and computational disclosure}
The relation data, exact verifier, independent branch-count implementation,
Sturm verifier, tests and logs accompany this article. Decisive calculations
use integer or rational arithmetic; floating-point roots are not used.
\bibliographystyle{amsplain}\bibliography{references}
\end{document}
